\clubpenalty100000000
\widowpenalty100000000
\documentclass[xetex, svgnames, 12pt]{scrartcl}

\usepackage{xeCJK}
\setCJKmainfont{BabelStone Han}
\usepackage{multicol}
\usepackage{enumitem}
\usepackage{setspace}
\usepackage[left=2.5cm, right=2.5cm, top=3cm, bottom=3cm]{geometry}
% XeLaTeX settings
\usepackage{xltxtra,polyglossia}
\setmainfont[Mapping=tex-text,Scale=1.0]{FreeSerif}
\setsansfont[Mapping=tex-text,Scale=1.0]{FreeSans}
\setmonofont{FreeMono}
\setmainlanguage{english}
\usepackage{makecell}
\usepackage{tabu}



\newfontfamily\tibetan{Tibetan Machine Uni}
\newfontfamily\sil{Charis SIL}
% Text styles
\newcommand\Bod[1]{\tibetan\textup{#1}}

\newcommand{\textunderset}[2]{\begin{tabular}[t]{@{}c@{}}#2\\[-0.3em]\scriptsize#1\end{tabular}}
\newcommand{\textoverset}[2]{\begin{tabular}[b]{@{}c@{}}\scriptsize#1\\[-0.38em]#2\end{tabular}}
\newcommand{\textoversetlarge}[2]{\begin{tabular}[b]{@{}c@{}}\scriptsize#1\\[-0.2em]#2\end{tabular}}



\newfontfamily\burm{Padauk}
\usepackage{scrlayer-scrpage}
\usepackage[
    colorlinks,
    urlcolor=black,
    linkcolor=black,
    citecolor=CornflowerBlue,
    filecolor=black,
    linktocpage=true,
    %dvipdfm
    ]{hyperref}

% more packages needed
\usepackage{tikz-cd}
\usepackage{amsfonts}
\usepackage{amsmath}
\usepackage{amssymb}
\usepackage{xcolor, color, colortbl, multirow}
% biblatex is used for bibtex rendering, style author-year
\usepackage[
    backend=bibtex,
    bibstyle=authoryear,
    giveninits=true,
    ibidtracker=strict,
    isbn=false,doi=false,url=false,
    labelnumber=true,
    loccittracker=strict,
    maxbibnames=10,
    maxcitenames=2,
    natbib=true,
    opcittracker=strict,
    sortcites=true,
    sorting=nyt,
    style=authoryear-ibid,
    alldates=terse,
    terseinits=false
    ]{biblatex}

\renewcommand*{\nameyeardelim}{\addspace}
\DeclareFieldFormat{postnote}{#1}
\DeclareFieldFormat{multipostnote}{#1}

\renewcommand\bibfont{\footnotesize}


% pagestyle settings
\pagestyle{scrheadings}
\ihead{%Hill
}
\chead{Chinese Dictionary}
\ohead{}
\ifoot{}
\cfoot{\pagemark}
\ofoot{}

\bibliography{asia}
\usepackage{setspace}
\usepackage{tabularx}
\usepackage{adjustbox}
\usepackage[]{tikz}
\graphicspath{{hard-character-images/}{images/}{../images/}}

\usetikzlibrary{calc,arrows,positioning,shapes}

\newcommand{\SET}[1]  {\ensuremath{\mathcal{#1}}}
\newcommand{\MAT}[1]  {\ensuremath{\boldsymbol{#1}}}
\newcommand{\VEC}[1]  {\ensuremath{\boldsymbol{#1}}}
\newcommand{\Video}{\SET{V}}
\newcommand{\video}{\VEC{f}}
\newcommand{\track}{x}
\newcommand{\Track}{\SET T}
\newcommand{\LMs}{\SET L}
\newcommand{\lm}{l}
\newcommand{\PosE}{\SET P}
\newcommand{\posE}{\VEC p}
\newcommand{\negE}{\VEC n}
\newcommand{\NegE}{\SET N}
\newcommand{\Occluded}{\SET O}
\newcommand{\occluded}{o}
%\definecolor{amazongreen}{cmyk}{0.83,0.24,0,0.12}
\definecolor{blue}{rgb}{0.0, 0.0, 0.55}
\definecolor{green}{rgb}{0.0, 0.5, 0.0}
\definecolor{red}{rgb}{0.75, 0.0, 0.2}
\definecolor{violet}{RGB}{153,153,255}
\definecolor{antiquewhite}{rgb}{0.98, 0.92, 0.84}
\definecolor{amazongreen}{RGB}{0,153,153}

%\newfontfamily\burm{Padauk}
\newfontfamily\tib{Tibetan Machine Uni}


%Define boxes

\tikzstyle{startstop} = [rectangle, rounded corners, minimum width=4.2cm, minimum height=1cm,text centered, draw=amazongreen, text=white, fill=red!30]
\tikzstyle{process} = [rectangle, minimum width=4.2cm, minimum height=1cm, text centered, draw=amazongreen, fill=orange!30, text=white]
%\tikzstyle{decision} = [diamond, minimum width=3cm, minimum height=1cm, text centered, draw=black, fill=green!30]
\tikzstyle{decision} = [rectangle, rounded corners,minimum width=4cm, minimum height=1cm, text centered, text width=3cm, text=white, fill=amazongreen]

% Draw a video
\newlength{\FSZ}
\newcommand{\drawvideo}[3]{% [0 0.25 0.5 0.75 1 1.25 1.5]
   \noindent\pgfmathsetlength{\FSZ}{\linewidth/#2}
   \begin{tikzpicture}[outer sep=0pt,inner sep=0pt,x=\FSZ,y=\FSZ]
   \draw[color=amazongreen!50!black] (0,0) node[outer sep=0pt,inner sep=0pt,text width=\linewidth,minimum height=0] (video) {\noindent#3};
   \path [fill=amazongreen!50!black,line width=0pt] 
     (video.north west) rectangle ([yshift=\FSZ] video.north east) 
    \foreach \x in {1,2,...,#2} {
      {[rounded corners=0.6] ($(video.north west)+(-0.7,0.8)+(\x,0)$) rectangle +(0.4,-0.6)}
    }
;
   \path [fill=amazongreen!50!black,line width=0pt] 
     ([yshift=-1\FSZ] video.south west) rectangle (video.south east) 
    \foreach \x in {1,2,...,#2} {
      {[rounded corners=0.6] ($(video.south west)+(-0.7,-0.2)+(\x,0)$) rectangle +(0.4,-0.6)}
    }
;
   \foreach \x in {1,...,#1} {
     \draw[color=amazongreen!50!black] ([xshift=\x\linewidth/#1] video.north west) -- ([xshift=\x\linewidth/#1] video.south west);
   }
   \foreach \x in {0,#1} {
     \draw[color=amazongreen!50!black] ([xshift=\x\linewidth/#1,yshift=1\FSZ] video.north west) -- ([xshift=\x\linewidth/#1,yshift=-1\FSZ] video.south west);
   }
   \end{tikzpicture}
}

\title{A dictionary of Chinese characteres in transliteration}
\author{Nathan W. Hill}
\begin{document}
% GENERATED FILE - DO NOT EDIT BY HAND.
\section*{Generated sample entries}

% Entry 18-01
\begin{multicols}{2}
\paragraph{\textoversetlarge{18-01}{\huge{可}}}
{\Large{qay}}, 
%ke3	
\textit{khaX};	%0001a
柯	%ke1	
{\large{{\textsuperscript{arb·}}
qay}}, 
\textit{ka};		%0001d
笴	%gan3	
{\large{\textsuperscript{bamb:}qay}},
\textit{kanX},		%0001e
%笴	%ge3	
\textit{kaX};		%0001e
河	%he2	
{\large{\textsuperscript{aq·}qay}},
\textit{ḫa};		%0001g
苛	%ke1	
{\large{\textsuperscript{herb:}qay}},
\textit{ḫa};		%0001i
呵	%he1	
{\large{\textsuperscript{or·}qay}},
\textit{ḫa};		%0001j
訶	%he1	
{\large{\textsuperscript{dic·}qay}},
\textit{ḫa};		%0001k
阿	%e1	
{\large{\textsuperscript{tum·}qay}},
\textit{{\sil{ꞏ}}a};		%0001m
㫊	%yi3	
\textit{{\sil{ꞏ}}aX},		%0001n
%㫊	%yi3	
\textit{{\sil{ꞏ}}ieX};		%0001n
\begin{itemize}[noitemsep]
\begin{spacing}{0.7} 

\item 
{\Large{何}}	%he2	
{\large{\textsuperscript{hom·}qay}}
= 
{\large{g\textoverset{a}{a}y}},
\textit{ḫa},		%0001f
%何	%he2	
\textit{ḫaX};		%0001f
荷	%he2	
{\large{{\textsuperscript{bamb:}}g\textoverset{a}{a}y}}, 
\textit{ḫa};		%0001o
    \item 
{\Large{哥}}	%ge1	
{\large{qay:qay}} = {\large{k\textoverset{a}{a}y}},
\textit{ka};		%0001p
歌	%ge1	
{\large{k\textoverset{a}{a}}y{\textsuperscript{·dehab}}}, 
\textit{ka};		%0001q
謌	%ge1	
{\large{{\textsuperscript{dic·}}k\textoverset{a}{a}y}}, 
\textit{ka};		%0001r
\item
{\Large{奇}}	%qi2	
{\large{\textsuperscript{magn:}qay}}
= {\large{q\textoverset{b}{a}y}},
\textit{gie},		%0001s
%奇	%ji1	
\textit{kie};		%0001s
琦	%qi2	
{\large{\textsuperscript{gem·}q\textoverset{b}{a}y}},
\textit{gie};		%0001t
騎	%qi2	
{\large{\textsuperscript{equ·}q\textoverset{b}{a}y}},
\textit{gie};		%0001u
錡	%qi2	
{\large{\textsuperscript{met·}q\textoverset{b}{a}y}},
\textit{gie},		%0001v
	%qi2	
\textit{gieX},		%0001v
%錡	%yi3	
\textit{ṅieX};		%0001v
寄
{\large{\textsuperscript{tect:}q\textoverset{b}{a}y}},
	%ji4	
\textit{kieH};		%0001x
掎	%ji3	
{\large{\textsuperscript{man·}q\textoverset{b}{a}y}},
\textit{kieX};		%0001y
畸	%ji1	
{\large{\textsuperscript{ag·}q\textoverset{b}{a}y}},
\textit{kie};		%0001z
羇	%ji1	
\textit{kie};		%0001a'
綺	%qi3	
{\large{\textsuperscript{ser·}q\textoverset{b}{a}y}},
\textit{khieX};		%0001b'
踦	%qi1	
{\large{\textsuperscript{ped·}q\textoverset{b}{a}y}},
\textit{khie},		%0001c'
%踦	%ji3	
\textit{kieX},		%0001c'
%踦	%yi3	
\textit{ṅieX};		%0001c'
攲	%qi1	
{\large{q\textoverset{b}{a}y\textsuperscript{·ram}}},
\textit{khie};		%0001d'
%攲	%ji1	
\textit{kie};		%0001d'
觭	%ji1	
{\large{\textsuperscript{corn·}q\textoverset{b}{a}y}},
\textit{khie};		%0001e'
倚	%yi3	
{\large{\textsuperscript{hom·}q\textoverset{b}{a}y}},
\textit{{\sil{ꞏ}}ieX};		%0001f'
椅	%yi1	
{\large{\textsuperscript{arb·}q\textoverset{b}{a}y}},
\textit{{\sil{ꞏ}}ie};		%0001g'
猗	%e1	
{\large{\textsuperscript{can·}q\textoverset{b}{a}y}},
\textit{{\sil{ꞏ}}aX},	%0001h'
%猗	%yi1	
\textit{{\sil{ꞏ}}ie},		%0001h'
%猗	%yi3	
\textit{{\sil{ꞏ}}ieX};		%0001h'
輢	%yi3	
{\large{\textsuperscript{vehic·}q\textoverset{b}{a}y}},
\textit{{\sil{ꞏ}}ieX};	%0001i'
陭	%yi1	
{\large{\textsuperscript{tum·}q\textoverset{b}{a}y}},
\textit{{\sil{ꞏ}}ie}.		%0001j'
\end{spacing}
\end{itemize}
\end{multicols}

% Entry 18-04
\begin{multicols}{2}
\paragraph{\textoversetlarge{18-04}{\huge{加}}}
{\large{kay}},
%加	%jia1	
\textit{kae};	%kae	%0015a
枷	%jia1	
{\large{{\textsuperscript{arb·}}kay}},
\textit{kae};	%kae	%0015c
%枷	%jia1	
\textit{kaeH};	%kaeH	%0015c
珈	%jia1	
{\large{{\textsuperscript{gem·}}kay}},
\textit{kae};	%kae	%0015d
駕	%jia4	
{\large{kay{\textsuperscript{:equ}}}},
\textit{kaeH};	%kaeH	%0015e
嘉	%jia1	
{\large{{\textsuperscript{tympan:}}kay}},
\textit{kae};	%kae	%0015g
哿	%ge3	
{\large{kay{\textsuperscript{:poss}}}},
\textit{kaX};	%kaX	%0015i
賀	%he4	
{\large{kay{\textsuperscript{:conch}}}},
\textit{ḫaH};	%haH	%0015j
\end{multicols}

% Entry 18-05
\begin{multicols}{2}
\paragraph{\textoversetlarge{18-05}{\huge{我}}}
%wo3	
{\Large{ṅay}}, 
\textit{ṅaX};		%0002a
俄	%e2	
{\large{{\textsuperscript{hom·}}
ṅay}}, 
\textit{ṅa};		%0002h
娥	%e2	
{\large{{\textsuperscript{fem·}}
ṅay}},
\textit{ṅa;}		%0002i
峨	%e2	
{\large{{\textsuperscript{mon·}}
ṅay}},
\textit{ṅa};		%0002k
睋	%e2	
{\large{{\textsuperscript{ocul·}}
ṅay}},
\textit{ṅa};		%0002l
莪	%e2	
{\large{{\textsuperscript{herb:}}ṅay}},
\textit{ṅa};		%0002m
誐	%e2	
{\large{{\textsuperscript{dic·}}ṅay}},
\textit{ṅa};		%0002n
餓	%e4	
{\large{{\textsuperscript{ed·}}ṅay}},
\textit{ṅaH};		%0002o
鵝	%e2	
{\large{ṅay{\textsuperscript{·av}}}},
\textit{ṅa};		%0002p
鵞	%e2	
{\large{ṅay{\textsuperscript{:av}}}},
\textit{ṅa};		%0002p
蛾	%e2	
{\large{{\textsuperscript{serp·}}ṅay}},
\textit{ṅa},		%0002q
	%yi3	
\textit{ṅieX}.\footnote{Were it not for the one reading \textit{ṅieX}, 我 {\large{ṅay}} could have been written {\large{ṅ\textoverset{a}{a}y}}.}		
\begin{itemize}[noitemsep]
\begin{spacing}{0.7} 
    \item 
{\Large{義}}	%yi4	
{\large{{\textsuperscript{ov:}}ṅay}} =
{\large{ṅ\textoverset{b}{a}y}},
\textit{ṅieH};		%0002r
儀	%yi2	
{\large{{\textsuperscript{hom·}}ṅ\textoverset{b}{a}y}},
\textit{ṅie};	%0002u
議	%yi4	
{\large{{\textsuperscript{dic·}}ṅ\textoverset{b}{a}y}},
\textit{ṅieH};		%0002v
蟻	%yi3	
{\large{{{\textsuperscript{serp·}}ṅ\textoverset{b}{a}y}}},
\textit{ṅieX};		%0002x
\begin{itemize}[noitemsep]
    \item 
{\Large{羲}}	%xi1	
{\large{{\textsuperscript{io·}}ṅ\textoverset{b}{a}y}} = 
{\large{ṅh\textoverset{b}{a}y}},
\textit{hie};		%0002y
犧	%xi1	
{\large{ṅh\textoverset{b}{a}y}},
\textit{hie}.		%0002z
\end{itemize}
\end{spacing}
\end{itemize}
\end{multicols}

% Entry 19-18
\begin{multicols}{2}
\begin{spacing}{0.7}
\paragraph{\textoversetlarge{19-18}{
\includegraphics[width=9mm]{U+26760.png}}}	
	%luo3	
{\large{r\textoverset{a}{o}y}},
\textit{lwa};	%lwa	%0014a
蠃	%luo3	
{\large{r\textoverset{a}{o}y{\textsuperscript{:serp}}}},
\textit{lwaX};	%lwaX	%0014b
羸	%lei2	
{\large{r\textoverset{a}{o}y{\textsuperscript{:ov}}}},
\textit{liwe};	%ljwe	%0014c
\end{spacing}

Baxter and Sagart reconstruct with -r.
\end{multicols}

% Entry 19-22
\begin{multicols}{2}
\begin{spacing}{0.7}
\paragraph{\textoversetlarge{19-22}{
\includegraphics[width=9mm]{U+27D2A.png}}}	
	%suo3	
{\large{s\textoverset{a}{o}y}},
\textit{swaX};	%swaX	%0013a
瑣	%suo3	
{\large{{\textsuperscript{gem·}}s\textoverset{a}{o}y}},
\textit{swaeX};	%swaeX	%0013b
%瑣	%suo3	
\textit{swaX};	%swaX	%0013b
\end{spacing}

\end{multicols}

% Entry 01-30
\begin{multicols}{2}
\begin{spacing}{0.7}
\paragraph{\textoversetlarge{01-30}{\huge{午}}}	
{\large{ṅa}},
%午	%wu3	
\textit{ṅuX};	%nguX	0060a
仵	%wu3	
{{\textsuperscript{hom·}}\large{ṅa}},
\textit{ṅuH},	%nguH	0060f
%仵	%wu3	
\textit{ṅuX};	%nguX	0060f
忤	%wu4	
{{\textsuperscript{cor·}}\large{ṅa}},
\textit{ṅuH};	%nguH	0060g
迕	%wu4	
{{\textsuperscript{amb·}}\large{ṅa}},
\textit{ṅuH};	%nguH	0060h
杵
%杵	%chu3	
{{\textsuperscript{arb·}}\large{ṅa}},
\textit{chiəX};	%tsyhoX	0086a
\begin{itemize}[noitemsep]
    \item {\Large{許}}	%xu3	
{\large{\textsuperscript{dic·}ṅa}} = {\large{ṅha}},
\textit{hiəX},	%xjoX	0060i
%許	%hu3	
\textit{huX};	%xuX	0060i
滸	%hu3	
{{\textsuperscript{aq·}}\large{ṅha}},
\textit{huX};	%xuX	0060k
\item {\Large{御}}	%yu4	
{\large{ṅa}} = {\large{ṅ\textoverset{b}{a}}},
\textit{ṅiəH};	%ngjoH	0060l
禦	%yu4	
 {\large{ṅ\textoverset{b}{a}{\textsuperscript{:spir}}}}, 
\textit{ṅiəX}.	%ngjoX	0060p
\end{itemize}
\end{spacing}

\begin{spacing}{0.7}
\end{spacing}
\end{multicols}

% Entry 01-38
\begin{multicols}{2}
\begin{spacing}{0.7}
\paragraph{\textoversetlarge{01-38}{\huge{者}}}	
%者	%zhe3	
{\large{ta}}, 
\textit{ciaeX};	%tsyaeX	0045a
赭	%zhe3	
{\large{{\textsuperscript{ruf·}}ta}},
\textit{ciaeX};	%tsyaeX	0045d
奢	%she1	
{\large{{\textsuperscript{magn:}}ta}}, 
\textit{śiae};	%syae	0045e
褚	%zhu3	
{\large{{\textsuperscript{vest·}}ta}},
\textit{ṭiəX};	%trjoX	0045g
帾	%du3	
{\large{{\textsuperscript{lint·}}ta}},
\textit{tuX};	%tuX	0045b'
睹	%du3	
{\large{{\textsuperscript{ocul·}}ta}},
\textit{tuX};	%tuX	0045c'
覩	%du3	
{\large{ta{\textsuperscript{·vid}}}},
\textit{tuX};	%tuX	0045d'
都	%du1	
{\large{ta{\textsuperscript{·urb}}}},
\textit{tu};	%tu	0045e'
楮	%chu3	
{\large{{\textsuperscript{arb·}}ta}},
\textit{ṭhiəX},	%trhjoX	0045i
%楮	%chu3	
\textit{tuX};	%tuX	0045i
屠	%tu2	
{\large{{\textsuperscript{ost·}}ta}},
\textit{du};	%du	0045i'
箸	%zhu4	
{\large{{\textsuperscript{bamb:}}ta}},
\textit{ḍiəH};	%drjoH	0045j
瘏	%tu2	
{\large{{\textsuperscript{infirm·}}ta}},
\textit{du};	%du	0045j'
渚	%zhu3	
{\large{{\textsuperscript{aq·}}ta}},
\textit{ciəX};	%tsyoX	0045k
煑	%zhu3	
{\large{ta{\textsuperscript{:ign}}}},
\textit{ciəX};	%tsyoX	0045l
煮	%zhu3	
{\large{ta{\textsuperscript{:ign}}}},
\textit{ciəX};	%tsyoX	0045m
䰞	%zhu3	
\footnote{SWJZ: 从鬲。⺸聲。}
\textit{ciəX};	%tsyoX	0045n
闍	%she2	
{\large{{\textsuperscript{portχ}}ta}},
\textit{jiae},	%dzyae	0045h'
%闍	%du1	
\textit{tu};	%tu	0045h'
{
\includegraphics[width=4.5mm]{U+26465.png}}
%⿱羽者
%zhu4	%Top: 羽, bottom: 者.
\textit{cəH};	%tsyoH	0045o
陼	%zhu3	
{\large{{\textsuperscript{urb·}}ta}},
\textit{ciəX};	%tsyoX	0045q
緒	%xu4	
{\large{{\textsuperscript{ser·}}ta}},
\textit{ziəX};	%zjoX	0045s
書	%shu1	
{\large{{\textsuperscript{pincern:}}ta}},
\textit{śiə};	%syo	0045t
暑	%shu3	
{\large{{\textsuperscript{sol:}}ta}},
\textit{śiəX};	%syoX	0045x
堵	%du3	
{\large{{\textsuperscript{terr·}}ta}},
\textit{tuX};	%tuX	0045y
\begin{itemize}[noitemsep]
    \item {\Large{豬}}	%zhu1	
{\large{{\textsuperscript{porc·}}ta}},
\textit{ṭiə};	%trjo	0045h
    瀦	%zhu1	
\textit{ṭiə};	%trjo	0045k'
    \item 
{\Large{著}}	%zhuo2	
{\large{{\textsuperscript{herb:}}ta}} = {\large{tak}},
\textit{ḍiak},	%drjak	0045n'
%著	%zhuo2	
\textit{ṭiak},	%trjak	0045n'
%著	%zhu4	
\textit{ṭiəH};	%trjoH	0045n'
躇	%chu2	
{\large{{\textsuperscript{ped·}}tak}},
\textit{ḍiə},	%drjo	0045o'
%躇	%chuo4	
\textit{ṭhiak};	%trhjak	0045o'
{
\includegraphics[width=4.5mm]{U+230C8.png}}
%[U+230C8]⿰著斤	%zhuo2	
{\large{{tak\textsuperscript{·libra}}}},
\textit{ṭiak};	%trjak	0045p'
    \item {\Large{諸}}	%zhu1	
{\large{{\textsuperscript{dic·}}ta}},
\textit{ciə};	%tsyo	0045p
    儲	%chu2	
\textit{ḍiə};	%drjo	0045l'
    \item {\Large{署}}	%shu3	
{\large{{\textsuperscript{nass:}}ta}} = {\large{d\textoverset{b}{a}s}},
\textit{jiəH};	%dzyoH	0045r
曙	%shu4	
{\large{{\textsuperscript{sol·}}d\textoverset{b}{a}s}},
\textit{jiəH}.	%dzyoH	0045m'
\end{itemize}
\end{spacing}

\begin{spacing}{0.7}
\end{spacing}
\end{multicols}

% Entry 01-43
\begin{multicols}{2}
\paragraph{\textoversetlarge{01-43}{\huge{予}}}	
{\large{la}},
%予	%yu2	
\textit{yiə},	%yo	0083a
%予	%yu3	
\textit{yiəX};	%yoX	0083a
{
\includegraphics[width=4mm]{U+24282.png}}	
{\large{\textsuperscript{amb·}la}},
\textit{ziə},
豫	 %yu4	
{\large{la\textsuperscript{·elephant}}},
\textit{yiəH};	%yoH	0083e
杼	%zhu4
{\large{\textsuperscript{arb·}la}},
\textit{ḍiəX},	%drjoX	0083f
%杼	%shu4	
\textit{jiəH},	%dzyoH	0083f
%杼	%shu4	
\textit{jiəX};	%dzyoX	0083f
抒	%shu1	
{\large{\textsuperscript{man·}la}},
\textit{jiəX};	%dzyoX	0083g
序	%xu4	
{\large{\textsuperscript{lax·}la}},
\textit{ziəX};	%zjoX	0083h
芧	%xu4	
{\large{\textsuperscript{herb·}la}},
\textit{ziəX};	%zjoX	0083i
紓	%shu1
{\large{\textsuperscript{ser·}la}},
\textit{śiə},	%syo	0083j
%紓	%shu1	
\textit{jiəX};	%dzyoX	0083j
舒	%shu1	
{\large{\textsuperscript{dom·}la}},
\textit{śiə};	%syo	0083k
野\footnote{I need to look into the relationship among 
野, {\includegraphics[width=4mm]{U+21428.png}}, and 埜.}	%ye3	
{\large{\textsuperscript{vic·}la}},
\textit{yiaeX};	%yaeX	0083l
{\includegraphics[width=4mm]{U+21428.png}}
{\large{\textsuperscript{xxx·}la}},
	%ye3	
%⿱㮊土
\textit{yiaeX};	%yaeX	0083m
埜	%ye3	
{\large{\textsuperscript{xxx·}la}},
\textit{yiaeX}.	%yaeX	0083n
\end{multicols}

% Entry 01-57
\begin{multicols}{2}
\paragraph{\textoversetlarge{01-57}{\huge{且}}}	
{\large{tsa}},
\textit{tshiaeX},
%且	%ju1	
\textit{tshiə},	%tshjo	0046a
%且	%ju1	
\textit{tsiə};	%tsjo	0046a
罝	%jie1	
{\large{{\textsuperscript{nass:}}tsa}},
\textit{tsiae};	%tsjae	0046h
柤	%zha1	
{\large{{\textsuperscript{arb·}}tsa}},
\textit{tṣae};	%tsrae	0046i
抯	%zha1	
{\large{{\textsuperscript{man·}}tsa}},
\textit{dziaeX},	%dzjaeX	0046j
%抯	%zha1	
\textit{tsiaeX},	%tsjaeX	0046j
%抯	%zha1	
\textit{tṣae};	%tsrae	0046j
蛆	%ju1	
{\large{{\textsuperscript{serp·}}tsa}},
\textit{tsiə};	%tsjo	0046m
岨	%qu1	
{\large{{\textsuperscript{mon·}}tsa}},
\textit{tshiə};	%tshjo	0046n
狙	%ju1	
{\large{{\textsuperscript{can·}}tsa}},
\textit{tshiə};	%tshjo	0046o
疽	%ju1	
{\large{{\textsuperscript{infirm·}}tsa}},
\textit{tshiə};	%tshjo	0046p
砠	%ju1	
{\large{{\textsuperscript{lap·}}tsa}},
\textit{tshiə};	%tshjo	0046q
雎	%ju1	
{\large{tsa{\textsuperscript{·passer}}}},
\textit{tshiə};	%tshjo	0046r
鴡	%ju1	
{\large{tsa{\textsuperscript{·av}}}},
\textit{tshiə};	%tshjo	0046s
苴	%ju1	
{\large{{\textsuperscript{herb:}}tsa}},
\textit{tshiə},	%tshjo	0046t
%苴	%ju1	
\textit{tsiə},	%tsjo	0046t
%苴	%zha3	
\textit{tṣaeX};	%tsraeX	0046t
咀	%ju3	
{\large{{\textsuperscript{or·}}tsa}},
\textit{dziəX};	%dzjoX	0046u
詛	%zu3	
{\large{{\textsuperscript{dic·}}tsa}},
\textit{tṣiəH};	%tsrjoH	0046x
阻	%zu3	
{\large{{\textsuperscript{tum·}}tsa}},
\textit{tṣiəX};	%tsrjoX	0046y
鉏
{\large{{\textsuperscript{met·}}tsa}},
	%zhu4	
\textit{dẓiəX};	%dzrjoX	0046a'
祖	%zu3	
{\large{{\textsuperscript{spirit·}}tsa}},
\textit{tsuX};	%tsuX	0046b'
組	%zu3	
{\large{{\textsuperscript{ser·}}tsa}},
\textit{tsuX};	%tsuX	0046e'
粗	%cu1	
{\large{{\textsuperscript{oryz·}}tsa}},
\textit{dzuX},	%dzuX	0046h'
%粗	%cu1	
\textit{tshu};	%tshu	0046h'
徂	%cu2	
{\large{{\textsuperscript{eo·}}tsa}},
\textit{dzu};	%dzu	0046i'
殂	%cu2	
{\large{{\textsuperscript{mal·}}tsa}},
\textit{dzu};	%dzu	0046j'
駔	%zu4	
{\large{{\textsuperscript{equ·}}tsa}},
\textit{dzuX},	%dzuX	0046m'
%駔	%zang3	
\textit{tsaṅX};	%tsangX	0046m'
楂	%zha1	
\textit{tṣae};	%tsrae	0046u'
\begin{itemize}[noitemsep]
    \item {\Large{租}}	%zu1	
{\large{{\textsuperscript{gran·}}tsa}} = {\large{tsa₂}},
\textit{tsu};	%tsu	0046d'
蒩	%zu1	
\textit{tsu};	%tsu	0046q'
    \item {\Large{沮}}	%ju3	
{\large{{\textsuperscript{aq·}}tsa}} = {\large{tsa₃}},
\textit{dziəX},	%dzjoX	0046k
%沮	%ju4	
\textit{tsiəH};	%tsjoH	0046k
菹	%zu1	
\textit{tṣiə};	%tsrjo	0046n'
\item  
{\Large{助}}	%zhu4	
{\large{tsa{\textsuperscript{·virt}}}}
 = {\large{tsa₄}},
\textit{dẓiəH};	%dzrjoH	0046z
耡	%chu2	
\textit{dẓiə},	%dzrjo	0046o'
%耡	%chu2	
\textit{dẓiəH};	%dzrjoH	0046o'
鋤	%chu2	
\textit{dẓiə};	%dzrjo	0046p'
\item {\Large{虘}}	%cuo2	
{\large{{\textsuperscript{tigr·}}tsa}}   = {\large{tsa₅}},
\textit{dzu};	%dzu	0046k'
%⿰黹虘	%chu3	
{
\includegraphics[width=4.2mm]{U+2A4D0.png}}
{\large{{\textsuperscript{acu·}}tsa}},
\textit{tṣhiəX};	%tsrhjoX	0046r'
樝	%zha1	
\textit{tṣae};	%tsrae	0046s'
{
\includegraphics[width=4.2mm]{U+20B6F.png}}
%⿰虘又	%zha1	
{\large{tsa{\textsuperscript{·dextr}}}},

\textit{tṣae};	%tsrae	0046v'
\end{itemize}


俎	%zu3	
\textit{tṣiəX};	%tsrjoX	0046v
not a XS character
\end{multicols}

% Entry 01-67
\begin{multicols}{2}
\begin{spacing}{0.7}
\paragraph{\textoversetlarge{01-67}{\huge{父}}}{\huge{(}}{
\includegraphics[width=6mm]{父.png}}{\huge{)}}
%父	%fu4	
{\large{pa}}, 
\textit{biuX},	%bjuX	0102a
%父	%fu3	
\textit{piuX};	%pjuX	0102a
釜	%fu3	
{\large{pa{\textsuperscript{:met}}}},
\textit{biuX};	%bjuX	0102f
斧	%fu3	
{\large{pa{\textsuperscript{:libra}}}},
\textit{piuX};	%pjuX	0102h
\begin{itemize}[noitemsep]
    \item 
{\Large{布}}({
\includegraphics[width=5mm]{布.png}})	%bu4	
{\large{pa{\textsuperscript{:lint}}}} = {\large{p\textoverset{a}{a}}},
\textit{puH};	%puH	0102j
怖	%bu4	
{{\textsuperscript{cor·}}p\textoverset{a}{a}}, 
\textit{phuH};	%phuH	0102m
    \item {\Large{甫}}({
\includegraphics[width=5mm]{甫.png}})	%bu4	
{\large{pa{\textsuperscript{:ager}}}} = {\large{pa₂}},
%fu3	
\textit{piuX};	%pjuX	0102n
脯	%fu3	
{\large{{\textsuperscript{carn·}}pa₂}},
\textit{piuX};	%pjuX	0102r
莆	%fu3	
{\large{{\textsuperscript{herb:}}pa₂}},
\textit{piuX};	%pjuX	0102s
黼	%fu3	
{\large{{\textsuperscript{acu·}}pa₂}},
\textit{piuX};	%pjuX	0102t
簠	%fu3	
{\large{{\textsuperscript{xxx·}}pa₂}},
\textit{piu}, %pju	0102u
%簠	%fu3	
\textit{piuH},	%pjuH	0102u
%簠	%fu3	
\textit{piuX};	%pjuX	0102u
輔	%fu3	
{\large{{\textsuperscript{vehic·}}pa₂}},
\textit{biuX};	%bjuX	0102v
鬴	%fu3	
{\large{{\textsuperscript{olla·}}pa₂}},
\textit{biuX};	%bjuX	0102y
圃	%pu3	
{\large{{\textsuperscript{cla·}}pa₂}},
\textit{puH},	%puH	0102z
%圃	%pu3	
\textit{puX};	%puX	0102z
補	%bu3	
{\large{{\textsuperscript{vest·}}pa₂}},
\textit{puX};	%puX	0102c'
逋	%bu1	
{\large{{\textsuperscript{amb·}}pa₂}},
\textit{pu};	%pu	0102d'
餔	%bu1	
{\large{{\textsuperscript{ed·}}pa₂}},
\textit{pu};	%pu	0102e'
痡	%pu1	
{\large{{\textsuperscript{infirm·}}pa₂}},
\textit{phiu},	%phju	0102g'
%痡	%pu1	
\textit{phu};	%phu	0102g'
鋪	%fu1	
{\large{{\textsuperscript{met·}}pa₂}},
\textit{phiu},	%phju	0102h'
%鋪	%pu1	
\textit{phu};	%phu	0102h'
哺	%bu3	
{\large{{\textsuperscript{or·}}pa₂}},
\textit{buH},	%buH	0102i'
酺	%pu2	
{\large{{\textsuperscript{vin·}}pa₂}},
\textit{bu},	%bu	0102k'
%酺	%pu2	
\textit{buH};	%buH	0102k'
匍	%pu2	
{\large{pa₂{\textsuperscript{·flect}}}},
\textit{bu};	%bu	0102l'
\begin{itemize}
\item 
{\Large{浦}}	
{\large{{\textsuperscript{aq·}}pa₂}} = {\large{p\textoverset{a}{a}₂}},
%pu3	
\textit{phuX};	%phuX	0102f'
蒲	%pu2	
{\large{{\textsuperscript{herb:}}p\textoverset{a}{a}₂}},
\textit{bu};	%bu	0102n'
\item 
{\Large{捕}}	%bu3	
{\large{{\textsuperscript{man·}}pa₂}} = {\large{p\textoverset{a}{a}₃}},
\textit{buH};	%buH	0102j'
蒱	%pu2	
{\large{\textsuperscript{herb:}p\textoverset{a}{a}₃}},
\textit{bu};	%bu	0102o'
\item {\Large{尃}}	%fu1	
{\large{pa₂{\textsuperscript{:digit}}}} = {\large{p\textoverset{b}{a}}},
\textit{phiu};	%phju	0102p'
傅	%fu4	
{\large{{\textsuperscript{hom·}}p\textoverset{b}{a}}}, 
\textit{biuH},	%bjuH	0102u'
%傅	%fu4	
\textit{piuH};	%pjuH	0102u'
榑	%fu2	
{\large{{\textsuperscript{arb·}}p\textoverset{b}{a}}}, 
\textit{biu};	%bju	0102v'
賻	%fu4	
{\large{{\textsuperscript{conch·}}p\textoverset{b}{a}}}, 
\textit{biuH}.	%bjuH	0102x'
\item 旉	%fu1	
{\large{pa₂{\textsuperscript{:quadr}}}} = {\large{p\textoverset{b}{a}₂}},
\textit{phiu};	%phju	0102q'
敷	%fu1	
{\large{p\textoverset{b}{a}₂}{\textsuperscript{·fer}}}, 
\textit{phiu};	%phju	0102t'
\end{itemize}
\end{itemize}
\end{spacing}
The \textit{xiéshēng} analysis of 簠 is not clear to me. 


\begin{spacing}{0.7}
\end{spacing}
\end{multicols}

\end{document}
